
\documentclass[11pt]{article}
\usepackage{amsmath,amssymb}
\usepackage{bm}
\usepackage[margin=1in]{geometry}
\usepackage{hyperref}

\begin{document}

\title{Incremental UB Parameterization for Stage A}
\author{}
\date{}
\maketitle

\noindent\textbf{Normative status.} This writeup is normative by reference from \texttt{docs/spec-db-core.md}. Any changes to the parameterization here MUST be mirrored in \texttt{docs/spec-db-core.md} and the corresponding DB-AT tests (e.g., DB-AT-026/027); exploratory alternatives SHALL be documented separately and marked non-normative.

\section{Conventions}

We fix a simple, consistent convention and use it throughout:

\begin{itemize}
  \item Miller indices: column vector $h = (h,k,l)^\top$.
  \item Reciprocal-space scattering vector in the lab frame:
        column vector $r^\ast \in \mathbb{R}^3$.
  \item Orientation / setting matrix $A^\ast \in \mathbb{R}^{3\times 3}$ defined by
        \begin{equation}
          r^\ast = A^\ast h.
        \end{equation}
  \item Unit cell parameters
        \begin{equation}
          \mathbf{c} = (a, b, c, \alpha, \beta, \gamma).
        \end{equation}
  \item A Busing--Levy type matrix $B(\mathbf{c})$ that depends \emph{only} on the cell
        and maps Miller indices into reciprocal-space vectors in a crystal-fixed frame.
  \item A rotation $U \in SO(3)$ that maps the crystal-fixed frame into the lab frame.
\end{itemize}

With these conventions, the setting matrix is
\begin{equation}
  A^\ast = U\,B(\mathbf{c}).
\end{equation}

\section{Baseline from DIALS/MOSFLM}

From dxtbx/DIALS/MOSFLM we already have:

\begin{itemize}
  \item A trusted initial orientation matrix $A^\ast_0$
        (e.g.\ from \verb|crystal.get_A()|).
  \item A trusted initial cell
        \begin{equation}
          \mathbf{c}_0 = (a_0, b_0, c_0, \alpha_0, \beta_0, \gamma_0).
        \end{equation}
\end{itemize}

We also know a specific formula for $B(\mathbf{c})$.
One common Busing--Levy style choice uses
\begin{align}
  \cos\alpha &= \cos(\alpha), &
  \cos\beta  &= \cos(\beta),  &
  \cos\gamma &= \cos(\gamma),
\end{align}
and cell volume
\begin{equation}
  V = a b c\,
      \sqrt{
        1 + 2\cos\alpha\cos\beta\cos\gamma
        - \cos^2\alpha - \cos^2\beta - \cos^2\gamma
      }.
\end{equation}
From these we derive reciprocal basis lengths and angles and assemble
$B(\mathbf{c})$ as a lower-triangular $3\times 3$ matrix (the standard
Busing--Levy recipe).

From the baseline cell we compute
\begin{equation}
  B_0 := B(\mathbf{c}_0),
\end{equation}
and define the baseline orientation matrix
\begin{equation}
  U_0 := A^\ast_0 B_0^{-1}.
\end{equation}
If our conventions match those of cctbx, $U_0$ should be numerically orthogonal:
\begin{equation}
  U_0^\top U_0 \approx I, \qquad \det U_0 \approx 1.
\end{equation}
Thus we have the ground-truth decomposition
\begin{equation}
  A^\ast_0 = U_0 B_0.
\end{equation}

\noindent
\emph{Important:} this decomposition is done once, at initialization.

\section{UB / Quaternion Round-Trip Test}

We want the UB/quaternion implementation to be internally consistent.
A strict round-trip test looks like this:

\begin{enumerate}
  \item Pick a cell $\mathbf{c}$ and build $B(\mathbf{c})$
        using the same function that will be used in production.
  \item Pick a rotation matrix $U \in SO(3)$.
        For testing, generate a random unit quaternion $q$ and convert it to
        $U(q)$.
  \item Form
        \begin{equation}
          A^\ast := U B(\mathbf{c}).
        \end{equation}
  \item Run this through the UB/quaternion machinery:
        \begin{itemize}
          \item Recompute $B' := B(\mathbf{c})$.
          \item Compute $U' := A^\ast {B'}^{-1}$.
          \item Convert $U' \to q' \to U''$ using
                \verb|matrix_to_quaternion| and \verb|quaternion_to_matrix|.
        \end{itemize}
  \item Check the errors
        \begin{equation}
          \|B' - B(\mathbf{c})\|,\quad
          \|U' - U\|,\quad
          \|U'' - U\|.
        \end{equation}
\end{enumerate}

We want relative errors at machine precision (e.g.\ $< 10^{-12}$ in double).
If these checks fail, there is a convention bug (transpose, using $B^{-1}$ where
$B$ was intended, wrong handedness, etc.).

This is what it means for the UB/A$^\ast$ implementation to ``pass the
round-trip test''.

\section{Incremental Parameterization in $U$ and Cell}

Instead of making the quaternion represent the \emph{absolute} orientation $U$,
we let it represent a \emph{misset} rotation $\Delta U$ on top of the baseline
$U_0$. Likewise, cell parameters are treated as increments around
$\mathbf{c}_0$.

\subsection{Orientation Increments with Quaternions}

Let the optimization variable be a unit quaternion
\begin{equation}
  q = (w, x, y, z)^\top \in \mathbb{R}^4, \qquad \|q\| = 1.
\end{equation}
We define the corresponding rotation matrix $R(q) \in SO(3)$ with the standard
formula
\begin{equation}
  R(q) =
  \begin{pmatrix}
    1 - 2(y^2 + z^2) & 2(xy - wz)     & 2(xz + wy) \\
    2(xy + wz)       & 1 - 2(x^2+z^2) & 2(yz - wx) \\
    2(xz - wy)       & 2(yz + wx)     & 1 - 2(x^2+y^2)
  \end{pmatrix}.
\end{equation}
The current orientation is then
\begin{equation}
  U(q) := R(q)\,U_0.
\end{equation}

Key points:
\begin{itemize}
  \item The identity orientation (no misset) is given by
        \begin{equation}
          q_{\text{id}} = (1,0,0,0)^\top,
          \qquad R(q_{\text{id}}) = I.
        \end{equation}
  \item At $q = q_{\text{id}}$ we have
        \begin{equation}
          U(q_{\text{id}}) = I U_0 = U_0.
        \end{equation}
\end{itemize}
Thus orientation increments are encoded as deviations of $q$ from
$(1,0,0,0)$; we never try to re-encode $U_0$ itself as a quaternion.
In code, $U_0$ is a fixed $3\times 3$ tensor and only $q$ is trainable.

\subsection{Cell Increments}

We already have a baseline cell $\mathbf{c}_0$.
Introduce a 6-vector of cell parameter increments
\begin{equation}
  \delta\mathbf{c} =
  (\delta a, \delta b, \delta c, \delta\alpha, \delta\beta, \delta\gamma)
\end{equation}
as trainable parameters (or log-scale versions for the lengths).
For an additive model, define the current cell
\begin{equation}
  \mathbf{c}(\delta\mathbf{c}) =
  (a_0 + \delta a,\;
   b_0 + \delta b,\;
   c_0 + \delta c,\;
   \alpha_0 + \delta\alpha,\;
   \beta_0 + \delta\beta,\;
   \gamma_0 + \delta\gamma).
\end{equation}
Alternatively, we can use multiplicative deltas for the lengths, e.g.
\begin{equation}
  a(\delta a) = a_0 e^{\delta a}, \quad \text{etc.}
\end{equation}
The current metric matrix is then defined as a pure function of the cell:
\begin{equation}
  B(\delta\mathbf{c}) := B\big(\mathbf{c}(\delta\mathbf{c})\big).
\end{equation}
We never derive $B$ from $A^\ast$; $B$ is always computed directly from
the cell.

\section{Forward Mapping: Parameters $\to$ UB $\to$ $A^\ast$ $\to$ Simulator}

Given refinement parameters $(q, \delta\mathbf{c})$ we build $A^\ast$ as follows:

\begin{enumerate}
  \item Normalize the quaternion:
        \begin{equation}
          \hat{q} = \frac{q}{\|q\|}.
        \end{equation}
  \item Compute the misset rotation and full orientation:
        \begin{equation}
          U = R(\hat{q})\,U_0.
        \end{equation}
  \item Compute the current cell and $B$ matrix:
        \begin{equation}
          \mathbf{c} = \mathbf{c}(\delta\mathbf{c}), \qquad
          B = B(\mathbf{c}).
        \end{equation}
  \item Build the setting matrix:
        \begin{equation}
          A^\ast(q,\delta\mathbf{c}) := U B
          = R(\hat{q})\,U_0\,B\big(\mathbf{c}(\delta\mathbf{c})\big).
        \end{equation}
\end{enumerate}

This $A^\ast(q,\delta\mathbf{c})$ is what we convert into
$(\texttt{mosflm\_a\_star}, \texttt{mosflm\_b\_star}, \texttt{mosflm\_c\_star})$
tuples and pass to nanoBragg / the simulator.

\subsection{Zero-Point Behaviour}

At the zero point of the parameters,
\begin{equation}
  q = q_{\text{id}} = (1,0,0,0), \qquad \delta\mathbf{c} = 0,
\end{equation}
we have
\begin{align}
  R(q_{\text{id}}) &= I, \\
  U(q_{\text{id}}) &= U_0, \\
  \mathbf{c}(0)    &= \mathbf{c}_0, \\
  B(0)             &= B(\mathbf{c}_0) = B_0,
\end{align}
and therefore
\begin{equation}
  A^\ast(q_{\text{id}}, 0) = U_0 B_0 = A^\ast_0.
\end{equation}
This equality holds by construction, not by an approximate quaternion
round-trip.

Along this forward path we never factor $A^\ast$ back into $(U,B)$.
We always go
\begin{equation}
  (q, \delta\mathbf{c})
  \;\longrightarrow\;
  (U,B)
  \;\longrightarrow\;
  A^\ast
  \;\longrightarrow\;
  \text{simulator}.
\end{equation}
All potentially lossy or convention-sensitive operations such as
$A^\ast_0 \to (U_0,B_0)$ and \verb|matrix_to_quaternion| are restricted to
\begin{itemize}
  \item a one-time decomposition $A^\ast_0 = U_0 B_0$ at initialization, and
  \item a pure forward map $q \to R(q)$ during refinement.
\end{itemize}

\section{Before vs After: Original vs Improved Formulation}

This section contrasts the original Stage-A quaternion/U-matrix formulation
with the improved incremental UB parameterization.

\subsection{Original Formulation (``Before'')}

Conceptual pipeline:

\begin{enumerate}
  \item Start from mapping $A^\ast_0$ (from DIALS/MOSFLM).
  \item Decompose into $(U_0, B_{\text{ideal}})$ such that
        $A^\ast_0 \approx U_0 B_{\text{ideal}}$ via some helper
        (e.g.\ \verb|derive_u_matrix_from_mosflm_a_star|).
  \item Encode the \emph{absolute} orientation as a quaternion:
        $q_0 = \texttt{matrix\_to\_quaternion}(U_0)$.
        The trainable parameter $q$ represents the full orientation, not a misset.
  \item Forward pass:
        \begin{itemize}
          \item Normalize $q$ to $\hat{q}$.
          \item Compute $U(q) = \texttt{quaternion\_to\_matrix}(\hat{q})$.
          \item Rebuild $A^\ast_{\text{recon}} = U(q) B_{\text{ideal}}$.
          \item Convert columns of $A^\ast_{\text{recon}}$ to
                $(\texttt{mosflm\_a\_star},\texttt{mosflm\_b\_star},\texttt{mosflm\_c\_star})$
                and pass to the simulator.
        \end{itemize}
  \item At the nominal ``zero'' point, the code implicitly expects
        $A^\ast_{\text{recon}}(q=q_0) \approx A^\ast_0$.
\end{enumerate}

In practice, this round trip often fails due to
\begin{itemize}
  \item inconsistent $B$ conventions (e.g.\ recomputing $B$ with a different formula),
  \item $A$ vs.\ $A^\ast$ or transpose mismatches,
  \item quaternion convention mismatches,
  \item float32/float64 and Python tuple conversions.
\end{itemize}
As a result, the orientation and hence $A^\ast$ are wrong even at
``zero'' parameters, causing large discrepancies in predicted peak positions
and catastrophic $\chi^2$ at initialization.

\subsection{Improved Formulation (``After'')}

The improved design treats the DIALS/dxtbx crystal state as authoritative and
uses \emph{incremental} parameters around this baseline.

Conceptual pipeline:

\begin{enumerate}
  \item Baseline decomposition (once):
        \begin{itemize}
          \item Get $A^\ast_0$ and $\mathbf{c}_0$ from dxtbx/DIALS.
          \item Compute $B_0 = B(\mathbf{c}_0)$ via a Busing--Levy-compatible
                formula.
          \item Define $U_0 = A^\ast_0 B_0^{-1}$ so that $A^\ast_0 = U_0 B_0$.
        \end{itemize}
  \item Define trainable increments:
        \begin{itemize}
          \item Orientation increment: unit quaternion $q$ with
                $q = q_{\text{id}}$ at initialization.
          \item Cell increments: $\delta\mathbf{c}$ with $\delta\mathbf{c}=0$
                at initialization.
        \end{itemize}
  \item Forward pass:
        \begin{itemize}
          \item Normalize $q$ to $\hat{q}$.
          \item Compute $U(q) = R(\hat{q}) U_0$.
          \item Compute $\mathbf{c} = \mathbf{c}(\delta\mathbf{c})$ and
                $B = B(\mathbf{c})$.
          \item Build $A^\ast(q,\delta\mathbf{c}) = U(q)\,B(\mathbf{c})$ and
                pass its columns as
                $(\texttt{mosflm\_a\_star},\texttt{mosflm\_b\_star},\texttt{mosflm\_c\_star})$
                to the simulator.
        \end{itemize}
\end{enumerate}

At zero parameters $(q=q_{\text{id}}, \delta\mathbf{c}=0)$ we have
$A^\ast(0) = A^\ast_0$ by construction, so the mapping and refinement zero
points are identical.

\paragraph{Advantages.}
\begin{itemize}
  \item The baseline crystal state $(U_0,B_0,A^\ast_0)$ comes directly from
        dxtbx/DIALS and is never re-encoded through a quaternion round trip.
  \item Orientation and cell are parameterized as incremental perturbations
        $(\Delta R, \Delta\text{cell})$ around the baseline, matching how
        DIALS and other frameworks conceptualize refinement.
  \item The simulator always consumes $A^\ast$ built as $U(\text{params}) B(\text{params})$
        in a single direction; no production code decomposes $A^\ast$ back into
        $(U,B)$, eliminating round-trip ambiguity inside the optimization loop.
\end{itemize}

\section{Summary}

The robust path can be summarized as
\begin{equation}
  A^\ast(q, \delta\mathbf{c})
  =
  R\!\left(\frac{q}{\|q\|}\right)
  U_0
  B(\mathbf{c}_0 + \delta\mathbf{c}),
\end{equation}
with $U_0$ and $\mathbf{c}_0$ extracted once from DIALS/MOSFLM.
All refinement happens in the incremental parameters $(q, \delta\mathbf{c})$;
we never re-factor $A^\ast$ back into $U$ and $B$ inside the optimization loop.

\end{document}
